%!TeX root=../main.tex

\chapter{Progettazione e implementazione del sistema}\label{ch:progettazione-e-implementazione-del-sistema}

\section{Linee guida progettuali}\label{sec:linee-guida-progettuali}
La progettazione di \textbf{Swift Contest} è stata guidata da principi di ingegneria del software moderni, volti a garantire un sistema robusto, sicuro e manutenibile, anche in un contesto di sviluppo rapido.
Le linee guida principali sono state:
\begin{itemize}
    \item \textbf{Separazione delle responsabilità (SoC)}: è stata applicata una netta separazione tra il livello di presentazione (UI), la logica di presentazione (ViewModel) e il livello di accesso ai dati (Model/Repository), fondamentale per la manutenibilità e scalabilità del codice.
    \item \textbf{Sicurezza come requisito primario}: è stato adottato un modello di accesso ai dati basato sul \textbf{Principio del Minimo Privilegio (Principle of Least Privilege)}, dove il client è considerato intrinsecamente non sicuro.
    L'accesso al database è negato per default e concesso selettivamente solo tramite un'API sicura.
    \item \textbf{Architettura Client-Server con backend BaaS}: è stata adottata una moderna architettura Client-Server.
    Il client (sviluppato in Flutter) gestisce la presentazione e la logica di UI, mentre il server è rappresentato da una piattaforma \textbf{Backend-as-a-Service (BaaS)} come Supabase.
    Questo approccio permette di delegare la gestione dell'infrastruttura (database, autenticazione, storage) e di distribuire la logica di business in modo flessibile tra il client e le funzioni serverless (RPC ed Edge Functions).
    \item \textbf{Sviluppo cross-platform}: la scelta di un framework come Flutter ha permesso di massimizzare il riuso del codice tra le piattaforme target (Android e Web), ottimizzando i tempi di sviluppo.
\end{itemize}

\begin{figure}[H]
    \centering
    \includegraphics[width=0.8\textwidth]{capitoli/figure/diagrams/client-server-architecture}
    \caption{Architettura Client-Server del sistema Swift Contest con backend BaaS.}
    \label{fig:client-server-architecture}
\end{figure}

\section{Architettura del frontend: il pattern MVVM con BLoC}
\label{sec:frontend-architecture}
Il frontend è stato sviluppato in Flutter seguendo i principi dell'architettura \textbf{MVVM (Model-View-ViewModel)}, implementata attraverso il pattern \textbf{BLoC (Business Logic Component)}.
Questa scelta permette di ottenere una netta separazione delle responsabilità, garantendo non solo \textbf{testabilità} e \textbf{manutenibilità}, ma anche la possibilità di \textbf{sviluppare la UI e la logica di business in parallelo}.
La struttura del codice sorgente in \texttt{lib/} riflette questa suddivisione:
\begin{itemize}
    \item \textbf{View (\texttt{/view})}: il livello di presentazione, composto dai Widget di Flutter che costituiscono l'interfaccia utente.
    Questo livello si limita a mostrare i dati ricevuti dal suo ViewModel (il BLoC) e a notificare gli eventi di interazione dell'utente (es.\ la pressione di un pulsante).
    \item \textbf{ViewModel (\texttt{/viewmodel}, implementato come \textbf{BLoC})}: il cuore della logica di presentazione.
    Ogni BLoC riceve eventi dalla View, esegue la logica di business (spesso invocando un metodo del Model/Repository), e produce uno stream di stati a cui la View reagisce per aggiornarsi.
    \item \textbf{Model (\texttt{/model})}: il livello di accesso ai dati.
    È composto dai \emph{Repository}, che fungono da unica interfaccia verso il backend, astraendo le chiamate a RPC ed Edge Functions.
    Include anche le \emph{Entities} (le rappresentazioni dei dati del dominio) e i \emph{Bundles}, ovvero oggetti aggregati che raggruppano dati correlati (es.\ un contest e i suoi partecipanti) in un'unica struttura, \textbf{ottimizzando il numero di chiamate di rete necessarie} per popolare una schermata.
\end{itemize}

\begin{figure}[H]
    \centering
    \includegraphics[width=0.8\textwidth]{capitoli/figure/diagrams/client-architecture}
    \caption{Architettura del Client MVVM con BLoC.}
    \label{fig:client-architecture}
\end{figure}

\subsection{Implementazione del ViewModel: il pattern BLoC}\label{subsec:implementazione-del-viewmodel:-il-pattern-bloc}
Per la gestione dello stato e l'implementazione del ViewModel, è stato adottato il pattern BLoC, utilizzando il pacchetto \texttt{flutter\_bloc}~\cite{bloc_docs}.
La scelta di BLoC è motivata dalla sua natura reattiva e basata su stream, che si integra perfettamente con il paradigma dichiarativo di Flutter.
Questo pattern impone un flusso di dati unidirezionale e prevedibile:
\begin{enumerate}
    \item La \textbf{View} notifica un \emph{Evento} (es. \texttt{LoginButtonPressed}).
    \item Il \textbf{BLoC} riceve l'evento, esegue la logica di business (es.\ chiama \texttt{authRepository.login()}) e, al termine, emette un nuovo \emph{Stato} (es. \texttt{LoginSuccess} o \texttt{LoginFailure}).
    \item La \textbf{View}, tramite un widget come \texttt{BlocBuilder}, si mette in ascolto dello stream di stati e si ridisegna per riflettere il nuovo stato.
\end{enumerate}
Questo approccio garantisce che la logica di business sia completamente disaccoppiata dalla UI, rendendola facilmente testabile in isolamento.
I BLoC sono stati organizzati in due categorie: \textbf{globali} (es. \texttt{AuthBloc}, \texttt{InboxBloc}) e \textbf{di pagina} (es. \texttt{JurorVotingProcedurePageBloc}).

\subsubsection{Gestione centralizzata degli errori nel repository}
Una delle responsabilità principali del livello \emph{Model} (e in particolare dei \emph{Repository}) è quella di agire come un robusto intermediario tra la logica di business e il backend.
Questo include non solo il recupero dei dati, ma anche una \textbf{gestione centralizzata e sicura degli errori}.

Le chiamate a un backend come Supabase possono fallire in molti modi: problemi di rete (\texttt{SocketException}), errori di autenticazione (\texttt{AuthException}), errori del database (\texttt{PostgrestException}), ecc.\ Esporre queste eccezioni direttamente al livello del BLoC o della UI sarebbe problematico per due motivi:
\begin{enumerate}
    \item Creerebbe un forte accoppiamento tra la logica di presentazione e le librerie specifiche del backend.
    \item Potrebbe esporre all'utente finale messaggi di errore tecnici che, oltre a essere incomprensibili, rappresentano un \textbf{rischio per la sicurezza} (information leakage), rivelando dettagli sull'implementazione del database o del backend.
\end{enumerate}

Per risolvere questo problema, è stata creata una funzione di ordine superiore, \texttt{handleBackendCall}, che ``avvolge'' ogni chiamata al backend.
Come mostrato nel listato seguente, questa funzione utilizza un blocco \texttt{try...catch} per intercettare le eccezioni specifiche di Supabase e mapparle in classi di \textbf{Failure} personalizzate e sicure.
Queste \emph{Failure} sono oggetti semplici che rappresentano un fallimento in modo astratto (es. \texttt{NetworkFailure}, \texttt{ServerFailure}), senza contenere dettagli tecnici.

Il risultato di questa mappatura viene quindi incapsulato utilizzando il tipo \textbf{\texttt{Either}}, fornito dal pacchetto di programmazione funzionale \texttt{fpdart}~\cite{fpdart_package}. \texttt{Either} è un tipo di dato che può contenere solo uno di due valori possibili: un valore di ``sinistra'' (\emph{Left}), che per convenzione rappresenta il fallimento, o un valore di ``destra'' (\emph{Right}), che rappresenta il successo.

Il tipo di ritorno \texttt{Either<Failure, T>} formalizza quindi questo contratto: la funzione restituirà o un oggetto di tipo \texttt{Failure} (a sinistra) o il dato atteso di tipo \texttt{T} (a destra).
Questo approccio elimina la necessità per il codice chiamante (il BLoC) di usare blocchi \texttt{try...catch}, permettendogli di gestire in modo pulito, dichiarativo e fortemente tipizzato sia il caso di successo che quello di fallimento, semplicemente ispezionando il contenuto di \texttt{Either}.

\begin{longlisting}
    \begin{minted}{dart}
        Future<Either<Failure, T>> handleBackendCall<T>(
          Future<Either<Failure, T>> Function() function,
        ) async {
            try {
                return await function();
            } on SocketException catch (e) {
                // Errore di rete (es. no internet).
                Logger.warning('Network error: $e');
                return Either.left(const NetworkFailure());
            } on PostgrestException catch (e) {
                // Errore del database PostgreSQL.
                Logger.warning('Postgrest error: $e');
                return Either.left(postgrestExceptionToFailure(e));
            }
            // ... altri catch specifici per le eccezioni di Supabase ...
            catch (e) {
                // Qualsiasi altro errore non previsto.
                Logger.error('Unknown error: $e');
                return Either.left(const Failure());
            }
        }
    \end{minted}
    \caption{Funzione helper per la gestione sicura delle chiamate al backend.}
    \label{lst:handle-backend-call}
\end{longlisting}

\subsubsection{Gestione dell'uguaglianza degli oggetti con Equatable}
Per ottimizzare le performance e garantire un comportamento prevedibile nella gestione dello stato, tutte le entità, i bundle e gli stati dei BLoC estendono la classe fornita dal pacchetto \texttt{equatable}~\cite{equatable_package}.

Questa scelta implementativa è cruciale per due motivi principali.

Il primo riguarda il corretto funzionamento del pattern BLoC. Di default, Dart confronta gli oggetti per riferimento, non per valore.
Ciò significa che due istanze di un oggetto, anche se contengono dati identici, vengono considerate diverse.
Senza \texttt{Equatable}, un BLoC potrebbe emettere un nuovo stato identico al precedente, causando una ricostruzione non necessaria e dispendiosa dell'interfaccia utente.
Utilizzando \texttt{Equatable}, invece, il confronto avviene basandosi sul valore degli attributi dell'oggetto.
Questo permette ai widget di Flutter, come \texttt{BlocBuilder}, di determinare in modo efficiente se uno stato è \textbf{realmente} cambiato prima di procedere con la ricostruzione, prevenendo refresh superflui della UI e migliorando significativamente la reattività dell'applicazione.

Il secondo motivo, altrettanto importante, riguarda la \emph{robustezza e la chiarezza della logica di business}.
Spesso, durante lo sviluppo, sorge la necessità di confrontare due oggetti per verificare se rappresentano la stessa entità concettuale, anche se sono istanze diverse in memoria.
Estendendo \texttt{Equatable}, si ottiene un'implementazione standard e priva di errori dell'uguaglianza basata sul valore, semplificando la logica di confronto all'interno dei BLoC e dei Repository e riducendo il rischio di bug sottili legati al confronto per riferimento.

\subsection{Navigazione con AutoRoute}\label{subsec:navigazione-con-autoroute}
La gestione della navigazione è stata affidata al pacchetto \texttt{auto\_route}~\cite{auto_route_package}, che genera un router fortemente tipizzato, eliminando l'uso di stringhe per le rotte.
Questo approccio trasforma potenziali errori di runtime (es.\ un errore di battitura nel nome di una rotta) in \textbf{errori di compilazione}, aumentando significativamente la robustezza dell'applicazione.
Un elemento chiave è l'\texttt{AuthGuard}, una classe che intercetta le richieste di navigazione e verifica lo stato di autenticazione dell'utente tramite l'\texttt{AuthBloc}, reindirizzandolo alla pagina di login se tenta di accedere a una rotta protetta.

\begin{longlisting}
    \begin{minted}{dart}
        // --- 1. Definizione della rotta protetta (in app_router.dart) ---
        AutoRoute(
            path: '/organizer-home',
            page: OrganizerHomeRoute.page,
            guards: [AuthGuard(authBloc: authBloc)], // Applica la guardia
        ),

        // ...

        // --- 2. Implementazione della guardia (in auth_guard.dart) ---
        class AuthGuard extends AutoRouteGuard {
            final AuthBloc authBloc;
            AuthGuard({required this.authBloc});

            @override
            void onNavigation(NavigationResolver resolver, StackRouter router) {
                // Controlla lo stato di autenticazione tramite l'AuthBloc
                final isAuthenticated = authBloc.state.authStatus.isAuthenticated;

                // Logica di base: se l'utente non è autenticato,
                // viene reindirizzato alla pagina di login.
                if (isAuthenticated) {
                    resolver.next(true); // Procede alla rotta protetta
                } else {
                    resolver.redirect(SignInRoute()); // Reindirizza al login
                }
            }
        }
    \end{minted}
    \caption{Estratto del sistema di routing protetto con ``AuthGuard''.}
    \label{lst:auth_guard}
\end{longlisting}

\subsection{Strategia di navigazione e passaggio dei dati}\label{subsec:strategia-di-navigazione-e-passaggio-dei-dati}
Una scelta architetturale fondamentale, dettata dalla necessità di una perfetta integrazione con il web, riguarda la modalità di passaggio dei dati tra le pagine.
A differenza di un'applicazione puramente mobile, dove è comune passare interi oggetti in memoria durante la navigazione, in un'applicazione web questo approccio è insostenibile: l'utente può ricaricare la pagina o accedere a un URL diretto, perdendo qualsiasi stato tenuto in memoria.

Per questo motivo, il sistema di navigazione di Swift Contest adotta una strategia rigorosa: \textbf{nelle rotte vengono passati solo identificativi (ID) primitivi} (es. \texttt{contestId}, \texttt{juryId}), mai oggetti complessi.
Quando una pagina viene caricata, il BLoC associato riceve l'ID dai parametri della rotta e, come prima azione, invoca il Repository per caricare dal backend tutti i dati necessari.

Questa scelta, sebbene aumenti il numero di chiamate al database, garantisce vantaggi in termini di robustezza e coerenza con il paradigma web:
\begin{itemize}
    \item \textbf{Ogni pagina è ricaricabile} in modo indipendente.
    \item \textbf{Ogni pagina ha un URL univoco} e condivisibile (deep linking).
    \item \textbf{Il comportamento dell'applicazione è prevedibile} e consistente tra la navigazione interna e l'accesso diretto a un URL\@.
\end{itemize}

\subsubsection{Validazione dell'input utente e gestione dei form}
Per garantire l'integrità dei dati e fornire un feedback immediato all'utente, è stata implementata una robusta strategia di validazione dell'input direttamente a livello client, prima che i dati vengano inviati al BLoC e, successivamente, al backend.

Invece di definire la logica di validazione all'interno dei singoli widget, \textbf{tutte le funzioni di validazione sono state centralizzate} in un unico file, \texttt{lib/utils/validators/validators.dart}.
Questo approccio, basato su funzioni pure, offre due vantaggi significativi:
\begin{itemize}
    \item \textbf{Riutilizzabilità}: La stessa funzione, ad esempio \texttt{emailValidator}, può essere riutilizzata in più punti dell'applicazione (es.\ pagina di login, registrazione, modifica profilo) senza duplicazione di codice.
    \item \textbf{Manutenibilità}: Se una regola di validazione deve essere modificata (es.\ la complessità minima di una password), la modifica viene effettuata in un unico punto, garantendo coerenza in tutta l'applicazione.
\end{itemize}

Il listato seguente mostra un esempio di un \emph{validator} e il suo utilizzo all'interno di un widget \texttt{TextFormField}.

\begin{longlisting}
    \begin{minted}{dart}
        // 1. Definizione della funzione di validazione pura
        String? passwordValidator(String? value) {
            if (value == null || value.isEmpty) {
                return 'Required';
            }
            if (value.length < 8) {
                return 'Must contain at least 8 characters';
            }
            // ... altre regole di validazione (maiuscole, numeri, etc.) ...

            return null; // La validazione ha successo
        }

        // 2. Utilizzo del validatore in un widget della UI
        TextFormField(
            validator: passwordValidator,
            autovalidateMode: AutovalidateMode.onUserInteraction,
            // ...
        );
    \end{minted}
    \caption{Esempio di validatore centralizzato e suo utilizzo in un widget.}
    \label{lst:validators}
\end{longlisting}

\subsection{Design system e theming}\label{subsec:design-system-e-theming}
Per garantire un'interfaccia utente coerente e facilmente manutenibile, è stato definito un design system centralizzato basato sul \textbf{Material Design 3}~\cite{material_design_docs}.
La scelta di \texttt{MaterialTheme} come base è stata strategica, in quanto è lo standard de facto per le applicazioni Android e garantisce un'esperienza utente nativa e riconoscibile.

Il cuore del sistema di theming di Material Design è il \textbf{\texttt{ColorScheme}}, un oggetto che definisce una palette di colori predefinita e semanticamente organizzata.
Invece di usare colori specifici i widget fanno riferimento a ruoli di colore come \texttt{primary} (per elementi principali), \texttt{secondary} (per elementi di supporto), \texttt{surface} (per gli sfondi), \texttt{onPrimary} (per il testo da mostrare sopra un colore primario), ecc\@.
Questo approccio permette di cambiare l'intera palette di colori dell'applicazione (ad esempio, passando dal tema chiaro a quello scuro) modificando unicamente l'oggetto \texttt{ColorScheme}, senza dover alterare il codice dei singoli widget.

Tuttavia, il \texttt{ColorScheme} standard non è sufficiente per coprire tutte le esigenze di un'applicazione complessa, che richiede colori semantici personalizzati (es.\ per stati di successo, errore, warning).

Per risolvere questo problema in modo pulito e scalabile, invece di definire colori statici sparsi nel codice, si è fatto ricorso a una funzionalità del linguaggio Dart: le \textbf{estensioni}.
Un'estensione (\emph{extension}) permette di \textbf{aggiungere nuove funzionalità a una classe esistente}, anche se non si ha accesso al suo codice sorgente.
In questo caso, è stata creata un'estensione sulla classe \texttt{ColorScheme} di Flutter per ``arricchirla'' con i colori semantici personalizzati richiesti dall'applicazione.

Come mostrato nel listato seguente, l'estensione \texttt{ColorSchemeX} aggiunge nuovi \emph{getter} (es. \texttt{success}, \texttt{warning}) alla classe \texttt{ColorScheme}.
Questi nuovi colori diventano così parte integrante del tema e possono essere richiamati in qualsiasi punto dell'applicazione in modo sicuro e coerente tramite \texttt{Theme.of(context).colorScheme.success}, ereditando automaticamente i cambiamenti tra tema chiaro e scuro.

\begin{longlisting}
    \begin{minted}{dart}
        // in lib/utils/themes/color_scheme_x.dart

        extension ColorSchemeX on ColorScheme {
            // Colore per indicare successo (es. notifiche positive)
            Color get success => const Color(0xFF28a745);
            Color get onSuccess => const Color(0xFFFFFFFF);

            // Colore per avvisi
            Color get warning => const Color(0xFFffc107);
            Color get onWarning => const Color(0xFF212529);

            // Altri colori personalizzati...
        }
    \end{minted}
    \caption{Definizione dell'estensione ``ColorSchemeX'' per aggiungere colori personalizzati al tema.}
    \label{lst:colorscheme_extension}
\end{longlisting}

\subsubsection{Gestione del feedback di caricamento: l'OverlayLoader}
Per fornire un feedback visivo all'utente durante le operazioni asincrone, è stato implementato un sistema di caricamento personalizzato, \texttt{OverlayLoader}.
Questo sistema utilizza l'Overlay di Flutter per mostrare un indicatore di progresso con uno sfondo semitrasparente e offuscato, permettendo all'utente di mantenere il contesto visivo della schermata sottostante.

L'implementazione si basa su due scelte di design chiave per garantire robustezza e facilità d'uso:
\begin{itemize}
    \item \textbf{Singleton Pattern}~\cite{gamma1995design}: La classe \texttt{OverlayLoader} è implementata come un singleton per garantire che esista una sola istanza del gestore del loader, prevenendo la visualizzazione di più overlay contemporaneamente.
    \item \textbf{Dart Extensions}: Per semplificare l'invocazione del loader, è stata creata un'estensione sulla classe \texttt{BuildContext}, che permette di usare una sintassi pulita e intuitiva come \texttt{context.showLoader()} e \texttt{context.hideLoader()}.
\end{itemize}

\paragraph{Gestione del ciclo di vita e prevenzione dei ``leak''}
Un aspetto cruciale per la robustezza di questo sistema è la gestione del suo ciclo di vita.
Per evitare che l'overlay del loader rimanga ``bloccato'' sullo schermo se una pagina viene chiusa (\emph{disposed}) durante un'operazione di caricamento (ad esempio, a causa di un errore o della navigazione dell'utente), è stata adottata una pratica difensiva: \textbf{il metodo \texttt{context.hideLoader()} viene sistematicamente invocato all'interno del metodo \texttt{dispose} di ogni pagina} che utilizza il loader.
Questo garantisce che, indipendentemente da come o perché la pagina viene distrutta, qualsiasi overlay di caricamento attivo venga sempre rimosso, prevenendo ``leak'' visuali e garantendo la stabilità dell'interfaccia utente.

\begin{longlisting}
    \begin{minted}{dart}
    // 1. Classe Singleton per la gestione dell'overlay
    class OverlayLoader {
        OverlayEntry? _entry;

        // Singleton pattern per garantire un'unica istanza
        OverlayLoader._();
        static final OverlayLoader _instance = OverlayLoader._();
        factory OverlayLoader() => _instance;

        void show(BuildContext context) {
            if (_entry != null) return; // Evita di mostrare più loader

            _entry = OverlayEntry(
              builder: (_) => const Positioned.fill(
                // ObscuredLoader è un widget custom con sfondo sfocato
                child: ObscuredLoader(),
              ),
            );

            // Inserisce l'entry nell'overlay radice per essere sempre in cima
            Overlay.of(context, rootOverlay: true).insert(_entry!);
        }

        void hide() {
            _entry?.remove();
            _entry = null;
        }
    }

    // 2. Estensione su BuildContext per un'API pulita e intuitiva
    extension OverlayLoaderX on BuildContext {
        void showLoader() => OverlayLoader().show(this);
        void hideLoader() => OverlayLoader().hide();
    }
    \end{minted}
    \caption{Implementazione del loader non-intrusivo tramite Overlay e Dart Extensions.}
    \label{lst:overlay_loader}
\end{longlisting}

\section{Progettazione del database}\label{sec:progettazione-del-database}
Il modello concettuale del dominio, descritto nel Capitolo 2, è stato tradotto in un modello logico e fisico implementato su un database \textbf{PostgreSQL}, gestito tramite la piattaforma Supabase.
La progettazione si è concentrata su tre pilastri: integrità dei dati, sicurezza e performance.

\begin{figure}[H]
    \centering
    \includegraphics[width=\textwidth]{capitoli/figure/diagrams/database}
    \caption{Diagramma Entità-Relazioni (ERD) dello schema del database.}
    \label{fig:database}
\end{figure}

\subsection{Schema fisico e tabelle principali}\label{subsec:schema-fisico-e-tabelle-principali}
Lo schema del database è composto da una serie di tabelle che implementano le entità del dominio.
Le tabelle cardine includono:
\begin{itemize}
    \item \texttt{profiles}: Questa tabella estende la tabella \texttt{auth.users} gestita internamente da Supabase.
    Poiché lo schema \texttt{auth} è protetto e non modificabile, per aggiungere attributi personalizzati (es.\ \texttt{full\_name}, \texttt{pref\_role}) è stata creata una tabella parallela nello schema \texttt{public}.
    La colonna \texttt{id} di \texttt{profiles} funge sia da chiave primaria che da chiave esterna verso \texttt{auth.users(id)}, stabilendo una relazione uno-a-uno.
    La creazione del profilo è automatizzata tramite un \textbf{trigger di database}, che inserisce un nuovo record in \texttt{profiles} ogni volta che un utente completa la registrazione.
    \item \texttt{contests}: memorizza le informazioni principali di un concorso.
    \item \texttt{participations} e \texttt{jurations}: implementano le classi di associazione che legano, rispettivamente, un utente a un contest (\emph{partecipazione}) e un utente a una giuria (\emph{appartenenza a una giuria}).
    \item \texttt{works}: contiene i dati dei lavori sottomessi dai partecipanti.
    \item \texttt{voting\_forms} e \texttt{voting\_form\_fields}: definiscono la struttura dei moduli di voto.
    \item \texttt{voting\_sessions}: rappresenta una sessione di voto.
    \item \texttt{voting\_session\_participants}, \texttt{voting\_session\_jurors}: rappresentano i membri coinvolti in una specifica sessione di voto, oltre a contenere gli snapshot dei dati per quella sessione.
    \item \texttt{voting\_form\_submissions} e \texttt{voting\_form\_submission\_values}: memorizzano i voti inviati dai giurati.
\end{itemize}

\subsection{Scelte di design chiave}\label{subsec:scelte-di-design-chiave}
Oltre alla semplice mappatura delle entità, sono state prese decisioni di design cruciali per garantire la robustezza del sistema.

\subsubsection{Integrità referenziale e strategia di cancellazione}
Per garantire la consistenza del database, è stato fatto un uso estensivo dei vincoli di chiave esterna.
Una decisione di design fondamentale è stata quella di adottare una strategia di \textbf{cancellazione fisica (hard delete)} piuttosto che una di \textbf{cancellazione logica (soft delete)}, principalmente per ottimizzare l'uso delle risorse (spazio su database e storage) e per semplificare la logica delle query di lettura.

Data questa scelta, la gestione della cancellazione dei dati correlati è diventata un aspetto critico, risolto con un approccio a più livelli.

\paragraph{Cancellazione a cascata tramite \texttt{ON DELETE CASCADE}}
Per le relazioni di composizione forte e dirette, dove l'entità ``figlia'' non ha dipendenze esterne come \emph{file nello Storage}, è stata utilizzata la politica \textbf{\texttt{ON DELETE CASCADE}}.
L'esempio più calzante è la relazione tra \texttt{voting\_forms} e \texttt{voting\_form\_fields}.
Poiché i campi di un form non hanno senso di esistere senza il form stesso, l'eliminazione di un record in \texttt{voting\_forms} provoca la cancellazione automatica e a cascata di tutti i record correlati in \texttt{voting\_form\_fields}.
Questo approccio delega interamente la logica di pulizia al database, garantendo massima efficienza e consistenza dei dati con il minimo sforzo implementativo.

\paragraph{Cancellazione tramite trigger per relazioni uno-a-uno}
In alcune relazioni uno-a-uno, la chiave esterna è stata posizionata nella tabella concettualmente ``principale'' per motivi di design.
L'esempio più significativo è la relazione tra \texttt{contests} e \texttt{places}: è la tabella \texttt{contests} a contenere il riferimento \texttt{place\_id} con un vincolo \texttt{UNIQUE}.

Questa configurazione, sebbene sia logicamente corretta, impedisce di usare la politica \texttt{ON DELETE CASCADE} per eliminare un \texttt{place} quando il \texttt{contest} associato viene cancellato.
Per risolvere questo problema e automatizzare la pulizia dei dati, è stato implementato un \textbf{trigger di database}.
Questo trigger si attiva \emph{dopo} (\texttt{AFTER}) l'eliminazione di un record dalla tabella \texttt{contests} ed esegue esplicitamente la cancellazione del record corrispondente nella tabella \texttt{places}, utilizzando l'ID del luogo del contest appena eliminato.
In questo modo, si garantisce l'integrità referenziale e si evita di lasciare record \texttt{places} orfani, delegando anche in questo caso la logica di pulizia al database.

\paragraph{Cancellazione di dati esterni tramite Edge Functions}
La cancellazione a cascata del database non può gestire dati esterni, come i file memorizzati in Supabase Storage.
Per questo motivo, le operazioni di cancellazione più complesse, come l'eliminazione di un utente o di un contest, sono gestite da \textbf{Edge Functions} dedicate (es.\ \texttt{user-delete-account}, \texttt{admin-delete-contest}). Queste funzioni agiscono da orchestratori: dopo aver eliminato i record dal database (che a loro volta attivano le cancellazioni a cascata e i trigger), invocano l'API di Supabase Storage per \textbf{eliminare tutti i file associati} (es.\ immagini del contest, file dei lavori), garantendo una pulizia completa e prevenendo la comparsa di file ``orfani''.
Questa logica non poteva essere implementata in una funzione RPC, poiché queste non hanno accesso all'API dello Storage.

\subsubsection{Snapshot dei dati tramite denormalizzazione controllata}
Per garantire l'immutabilità e l'auditabilità del processo di voto, è stato implementato un \textbf{meccanismo di snapshot dei dati}, una delle scelte architetturali più importanti del backend.

Quando un organizzatore avvia una \emph{VotingSession} tramite la funzione RPC \texttt{organizer\_start\_voting\_session}, il sistema non si limita a collegare le entità esistenti, ma esegue una forma di \textbf{denormalizzazione controllata}, \emph{duplicando intenzionalmente} alcuni dati per disaccoppiare il processo di voto dallo stato corrente del contest.
Per ogni partecipante e giurato coinvolto, vengono create nuove righe in tabelle di sessione dedicate (es.\ \texttt{voting\_session\_participants}, \texttt{voting\_session\_jurors}).

Queste nuove righe contengono non solo le chiavi esterne ai record originali, ma anche una copia statica degli attributi più importanti al momento dell'avvio.
Come mostrato nel listato seguente, la tabella \texttt{voting\_session\_participants} include, oltre a \texttt{participation\_id}, anche i campi \texttt{participant\_full\_name}, \texttt{work\_name}, ecc.

Una scelta di design cruciale è l'uso della clausola \textbf{\texttt{ON DELETE SET NULL}} sulla chiave esterna \texttt{participation\_id}.
Questo significa che se il partecipante originale dovesse essere eliminato dal sistema, il record dello snapshot \emph{non} verrebbe cancellato, ma manterrebbe solo il riferimento a \texttt{NULL}.
Ciò dimostra che la vera ``fonte di verità'' per la sessione di voto sono gli attributi di snapshot, garantendo che i voti espressi rimangano validi e auditabili anche se le entità originali cessano di esistere.

Lo stesso principio si applica a \texttt{voting\_session\_jurors}, dove \texttt{juror\_id} è fondamentale per tracciare se un giurato ha già votato, ma può essere impostato a \texttt{NULL} in caso di cancellazione dell'account senza invalidare il voto già espresso.

Questa tecnica garantisce:
\begin{enumerate}
    \item \textbf{Integrità della sessione}: se un utente cambia nome, il nome ``congelato'' nello snapshot rimane invariato.
    \item \textbf{Auditabilità e resilienza}: i voti rimangono validi e associati ai dati di snapshot anche se i record originali vengono eliminati.
    \item \textbf{Performance ottimizzate}: i dati di snapshot vengono letti direttamente, evitando query con join.
\end{enumerate}

% Estratto del DDL SQL che mostra la definizione degli attributi di snapshot.
\begin{longlisting}
    \begin{minted}{postgresql}
    -- Definizione della tabella `voting_session_participants`

    CREATE TABLE public.voting_session_participants (
    id uuid PRIMARY KEY DEFAULT extensions.uuid_generate_v4(),
    voting_session_id uuid NOT NULL REFERENCES public.voting_sessions(id) ON DELETE CASCADE,

    -- Riferimento opzionale al record di partecipazione originale.
    -- ON DELETE SET NULL: se il partecipante viene rimosso dal contest,
    -- lo snapshot sopravvive, garantendo l'integrità dei voti.
    participation_id uuid REFERENCES public.participations(id) ON DELETE SET NULL,

    order_index int NOT NULL,

    -- --- Attributi di Snapshot ---
    -- Dati "congelati" al momento dell'avvio della sessione.
    participant_full_name varchar(70) NOT NULL,
    work_name varchar(50) NOT NULL,
    work_description varchar(1000) NOT NULL,
    work_images_paths text[] NOT NULL,

    UNIQUE (voting_session_id, participation_id),
    UNIQUE (voting_session_id, order_index)
    );

    -- Definizione della tabella `voting_session_jurors`

    CREATE TABLE public.voting_session_jurors (
    id uuid PRIMARY KEY DEFAULT extensions.uuid_generate_v4(),
    voting_session_id uuid NOT NULL REFERENCES public.voting_sessions(id) ON DELETE CASCADE,
    voting_session_jury_id uuid NOT NULL REFERENCES public.voting_session_juries(id) ON DELETE CASCADE,

    -- Riferimento opzionale al giurato.
    -- ON DELETE SET NULL: se l'account del giurato viene eliminato,
    -- il suo voto rimane registrato ma anonimizzato.
    juror_id uuid REFERENCES public.profiles(id) ON DELETE SET NULL,

    has_submitted bool NOT NULL DEFAULT false,

    -- --- Attributo di Snapshot ---
    juror_full_name varchar(70) NOT NULL,

    UNIQUE (voting_session_jury_id, juror_id)
    );
    \end{minted}
    \caption{Definizione delle tabelle di sessione con attributi di snapshot e gestione della cancellazione.}
    \label{lst:snapshot_attributes}
\end{longlisting}

\subsection{Implicazioni del design del form di voto sull'aggregazione dei risultati}\label{subsec:implicazioni-del-design-del-form-di-voto-sull'aggregazione-dei-risultati}
Una scelta di design fondamentale, visibile nello schema del database, è che \textbf{ogni giuria è associata a un unico e specifico form di voto}.
Questa flessibilità permette all'organizzatore di creare criteri di valutazione diversi per giurie diverse (es.\ una giuria tecnica con campi slider complessi e una giuria popolare con un singolo campo di preferenza).

Tuttavia, questa decisione architetturale ha un'implicazione diretta e importante su come i risultati possono essere processati e visualizzati, influenzando sia il backend che il frontend: \textbf{il sistema non può aggregare automaticamente i punteggi tra giurie diverse}, poiché i loro form di voto (e quindi i loro criteri e punteggi) non sono direttamente comparabili.

Di conseguenza, tutte le operazioni di analisi dei risultati, sia a livello di logica di backend (RPC) che di interfaccia utente (frontend), sono state progettate per operare a livello della singola giuria:
\begin{itemize}
    \item La \textbf{generazione della classifica} avviene per ogni singola giuria.
    \item L'\textbf{esportazione dei risultati} in formato Excel è granulare e produce un file per ciascuna giuria.
\end{itemize}
Se l'organizzatore desidera creare una classifica ``globale'' che combini i risultati di più giurie, deve farlo esternamente, utilizzando i dati esportati.
Questa è una limitazione consapevole del progetto, un compromesso accettato per garantire la massima flessibilità nella configurazione dei criteri di valutazione.

Per mantenere questa flessibilità anche nella fase finale, il sistema non impone una classifica globale, ma permette all'organizzatore di \textbf{pubblicare qualsiasi classifica in formato PDF}.
Questa scelta di design gli conferisce la \textbf{massima libertà}: può decidere di pubblicare la classifica di una singola giuria, oppure può elaborare esternamente una classifica globale personalizzata e caricarla come file PDF. In entrambi i casi, una volta pubblicata, la classifica diventa disponibile per il download a tutti i partecipanti e giurati del contest all'interno dell'applicazione.

\section{Architettura del backend}\label{sec:architettura-del-backend}
Il backend è costruito interamente su Supabase e la sua architettura è definita da un'API sicura composta da un insieme di Funzioni RPC ed Edge Functions.

\subsection{Funzioni RPC (Remote Procedure Call)}\label{subsec:funzioni-rpc-(remote-procedure-call)}
Le funzioni RPC, scritte in PL/pgSQL, sono il principale punto di ingresso per le operazioni sul database.
La scelta di utilizzare le RPC come strato di accesso ai dati primario, invece di eseguire query dirette dal client, è stata dettata da tre motivi fondamentali:

\begin{enumerate}
    \item \textbf{Sicurezza e Controllo}: Sono tutte dichiarate con \texttt{SECURITY DEFINER}, il che significa che vengono eseguite con i privilegi del loro creatore (un ruolo amministrativo) e non con quelli dell'utente chiamante.
    Al loro interno, contengono logica di validazione e controllo dei permessi, tipicamente verificando l'identità dell'utente tramite \texttt{auth.uid()}.
    Questa scelta è il fulcro dell'approccio di \textbf{accesso mediato da API}: nessun utente può effettuare operazioni dirette sul database, ma deve passare attraverso questo strato di controllo.

    \item \textbf{Supporto per le Transazioni}: Una limitazione chiave della libreria client \texttt{supabase-flutter} è la \textbf{mancanza di supporto per le transazioni multi-statement}.
    Operazioni complesse che richiedono l'inserimento o l'aggiornamento di record su più tabelle in modo atomico (es.\ la creazione di un contest e del suo luogo associato) non sarebbero sicure se eseguite tramite chiamate multiple dal client.
    Implementando questa logica all'interno di una singola funzione RPC, si sfrutta la capacità nativa di PostgreSQL di eseguire tutte le operazioni in un'unica transazione, garantendo che l'operazione abbia successo o fallisca nella sua interezza, senza lasciare il database in uno stato inconsistente.

    \item \textbf{Performance e Astrazione}: Le RPC permettono di incapsulare logica complessa e query con join multipli lato server, riducendo la quantità di dati trasferiti e il numero di round-trip tra client e server.
\end{enumerate}

% Elenco delle principali Funzioni RPC del sistema.
\begin{longtable}{p{0.3\textwidth} p{0.4\textwidth} p{0.25\textwidth}}
    \toprule
    \textbf{Nome Funzione} & \textbf{Scopo Principale} & \textbf{Controllo di Sicurezza} \\
    \midrule
    \endfirsthead
    \toprule
    \textbf{Nome Funzione} & \textbf{Scopo Principale} & \textbf{Controllo di Sicurezza} \\
    \midrule
    \endhead

    \texttt{\seqsplit{auth\_get\_account\_bundle}} &
    Ottiene i dati completi dell'utente autenticato.
    &
    \texttt{auth.uid() = id} \\
    \midrule

    \texttt{\seqsplit{organizer\_create\_contest}} &
    Crea un nuovo contest e il luogo associato.
    &
    Verifica che l'utente sia autenticato. \\
    \midrule

    \texttt{\seqsplit{organizer\_start\_voting\_session}} &
    Avvia una sessione di voto, creando lo snapshot dei dati.
    &
    L'utente deve essere l'organizzatore del contest. \\
    \midrule

    \texttt{\seqsplit{participant\_join\_contest}} &
    Associa un partecipante a un contest tramite un token di invito.
    &
    Il token deve essere valido e non utilizzato. \\
    \midrule

    \texttt{\seqsplit{juror\_submit\_votes}} &
    Salva i voti di un giurato per una sessione.
    &
    La sessione deve essere `live' e il giurato non deve aver già votato. \\
    \midrule

    \texttt{\seqsplit{juror\_access\_voting\_as\_simple\_juror}} &
    Fornisce l'accesso a un giurato ospite tramite il token della giuria.
    &
    Il token della giuria deve essere valido. \\

    \bottomrule
    \caption{Elenco esemplificativo delle Funzioni RPC del sistema.}
    \label{tab:rpc_functions}
\end{longtable}

% Inserire qui un estratto significativo dalla funzione PL/pgSQL `juror_submit_votes`, mostrando la validazione di un campo slider prima dell'inserimento.
\begin{longlisting}
    \begin{minted}{postgresql}
    -- STEP 2.5: VALIDATE VOTE VALUES
    -- Controlla in modo efficiente se qualche valore di tipo 'slider' nel payload JSON
    -- è al di fuori del suo range min/max definito o non è un numero valido.
    IF EXISTS (
      -- Utilizza una Common Table Expression (CTE) per "spacchettare" il payload JSON
      -- in una tabella temporanea di righe e colonne.
      WITH unpacked_votes AS (
        SELECT
          (value->>'voting_form_field_id')::uuid as field_id,
          (value->>'value')::text as value
        FROM jsonb_array_elements(p_votes_payload)
      )
      -- Seleziona le righe che violano le regole di validazione.
      SELECT 1
      FROM unpacked_votes uv
      -- Unisce i voti spacchettati con la tabella dei campi del form
      -- per ottenere i metadati di ogni campo (tipo, min, max).
      JOIN public.voting_form_fields vff ON vff.id = uv.field_id
      WHERE
        -- Applica il controllo solo ai campi di tipo 'slider'.
        vff.type = 'slider'
        AND (
          -- Condizione 1: Il valore non è una stringa di un numero intero.
          NOT (uv.value ~ '^-?\d+$')
          -- Condizione 2: Oppure, se è un intero, è fuori dai limiti.
          OR uv.value::integer < vff.slider_min_value
          OR uv.value::integer > vff.slider_max_value
        )
    ) THEN
      -- Se anche solo una riga viola le condizioni, solleva un'eccezione
      -- e interrompe la transazione, impedendo l'inserimento dei dati.
      RAISE EXCEPTION 'Invalid vote value: A slider value is outside its allowed range or is not a valid number.';
    END IF;

    -- Se la validazione passa, la funzione procede con l'inserimento dei voti...
    \end{minted}
    \caption{Estratto dalla funzione ``juror\_submit\_votes'' che mostra la validazione server-side dei valori degli slider.}
    \label{lst:juror_submit_votes_validation}
\end{longlisting}

\subsection{Edge Functions}\label{subsec:edge-functions}
Le Edge Functions, scritte in TypeScript ed eseguite su Deno, gestiscono operazioni complesse, transazioni multi-step e integrazioni con servizi esterni che non sarebbero possibili o sicure da eseguire in una semplice RPC\@.

% Elenco delle principali Edge Functions del sistema.
\begin{longtable}{p{0.25\textwidth} p{0.4\textwidth} p{0.25\textwidth}}
    \toprule
    \textbf{Nome Funzione} & \textbf{Scopo} & \textbf{Autorizzazione} \\
    \midrule
    \endfirsthead
    \toprule
    \textbf{Nome Funzione} & \textbf{Scopo} & \textbf{Autorizzazione} \\
    \midrule
    \endhead

    \texttt{\seqsplit{organizer-invite-juror}} & Invia un'email di invito a un giurato.
    & JWT + Ownership del contest. \\
    \midrule

    \texttt{\seqsplit{participant-submit-work}} & Gestisce il caricamento del file \texttt{.zip} del lavoro e delle immagini di anteprima.
    & JWT + Appartenenza al contest. \\
    \midrule

    \texttt{\seqsplit{user-delete-account}} & Elimina un utente e tutti i suoi dati e file associati.
    & JWT (invocata dall'utente stesso). \\
    \midrule

    \texttt{\seqsplit{admin-delete-account}} & Versione amministrativa dell'eliminazione account.
    & \textbf{X-API-Key} (per Retool). \\
    \midrule

    \texttt{\seqsplit{google-places-search-suggestions}} & Funge da proxy sicuro per l'API di Google Places.
    & JWT. \\
    \midrule

    \texttt{\seqsplit{get-latest-apk}} & Proxy sicuro per il download dell'APK da un repository GitHub privato.
    & JWT. \\

    \bottomrule
    \caption{Elenco esemplificativo delle Edge Functions del sistema.}
    \label{tab:edge_functions}
\end{longtable}

% Inserire qui un estratto di codice TypeScript dalla Edge Function `organizer-invite-juror` che mostra l'invocazione dell'API di Resend.
\begin{longlisting}
    \begin{minted}{typescript}
    // Import e inizializzazione del client di Resend.
    // La chiave API viene letta in modo sicuro dai segreti di Supabase.
    import { Resend } from 'npm:resend';
    const resend = new Resend(Deno.env.get('RESEND_API_KEY'));

    // ... (all'interno della funzione Deno.serve, dopo aver validato l'input
    // e aver creato il record di invito nel database)

    try {
        // ...
        // Passaggio 5: Inserisce il nuovo invito nel DB e ottiene il token generato.
        const { data: newInvitation, error: insertError } = await supabaseAdmin
            .from('juror_invitations')
            .insert({ contest_id, jury_id, email })
            .select()
            .single();

        if (insertError) throw insertError;

        // Passaggio 6: Recupera i nomi del contest e della giuria per personalizzare l'email.
        const { data: contest } = await supabaseAdmin.from('contests').select('name').eq('id', contest_id).single();
        const { data: jury } = await supabaseAdmin.from('juries').select('name').eq('id', jury_id).single();

        // Passaggio 7: Invia l'email di invito utilizzando il client di Resend.
        // L'operazione è asincrona.
        await resend.emails.send({
            from: "Swift Contest <noreply@swiftcontest.com>",
            to: [email],
            subject: `Invitation to join the jury for "${contest.name}"`,
            // Il corpo dell'email è un template HTML che include i dati dinamici.
            html: `
            <div>
            <h1>Invitation to Swift Contest</h1>
            <p>You have received an invitation to join the "<strong>${jury.name}</strong>" jury for the contest "<strong>${contest.name}</strong>".</p>
            <p>Use this token in the app to accept:<br>
            <strong>${newInvitation.token}</strong>
            </p>
            </div>
            `,
        });

        // Passaggio 8: Se l'invio ha successo, restituisce l'invito creato al client.
        return new Response(JSON.stringify(newInvitation), {
            headers: { ...corsHeaders, "Content-Type": "application/json" },
            status: 201, // Created
        });

    } catch (error) {
        // Se l'invio dell'email o qualsiasi operazione precedente fallisce,
        // l'errore viene catturato e registrato, e viene restituita una risposta di errore.
        console.error("Error sending invitation:", error.message);
        return new Response(JSON.stringify({ error: error.message }), {
            headers: { ...corsHeaders, "Content-Type": "application/json" },
            status: 500,
        });
    }
    \end{minted}
    \caption{Estratto dalla Edge Function ``organizer-invite-juror'' che mostra l'invocazione dell'API di Resend.}
    \label{lst:resend_invocation}
\end{longlisting}

\section{Architettura di Sicurezza: Principio del Minimo Privilegio e Accesso Mediato}\label{sec:architettura-di-sicurezza:-principio-del-minimo-privilegio-e-accesso-mediato}
La sicurezza è un requisito trasversale che ha guidato l'intera progettazione del backend.
È stato adottato un rigoroso modello basato sul \textbf{Principio del Minimo Privilegio (PoLP)}, il cui assunto fondamentale è negare l'accesso per default e concederlo solo in modo esplicito e controllato.

\subsection{Strategia di Implementazione: Centralizzazione della Logica di Autorizzazione}\label{subsec:strategia-di-implementazione:-centralizzazione-della-logica-di-autorizzazione}
Per implementare il Principio del Minimo Privilegio, si sarebbero potute seguire due strategie principali:
\begin{enumerate}
    \item Definire policy di \textbf{Row Level Security (RLS)}~\cite{postgresql_docs} granulari per ogni tabella e per ogni operazione (SELECT, INSERT, UPDATE, DELETE).
    \item Centralizzare tutti i controlli di autorizzazione all'interno di un'API di accesso ai dati (RPC e Edge Functions).
\end{enumerate}

Sebbene un approccio ibrido con un doppio livello di controllo (sia RLS che controlli nell'API) rappresenti la massima sicurezza teorica, per questo progetto è stata adottata la \textbf{seconda strategia}.
La logica di autorizzazione è stata interamente centralizzata all'interno delle Funzioni RPC e delle Edge Functions.
Questa scelta è stata un \textbf{compromesso consapevole}, motivato dalla necessità di ottimizzare i tempi di sviluppo e la manutenibilità: gestire una logica di autorizzazione complessa in un unico punto (il codice della funzione) è risultato più rapido e meno prono a errori rispetto a dover mantenere sincronizzate decine di policy RLS scritte in SQL.

\subsection{Il Ruolo della RLS e dell'API come "Gatekeeper"}\label{subsec:il-ruolo-della-rls-e-dell'api-come-gatekeeper}
In questo modello, la RLS gioca comunque un ruolo fondamentale, ma più semplice: \textbf{agisce come un "firewall" di default}.
Su tutte le tabelle è stata abilitata la RLS con una policy di \texttt{DENY} implicita.
Questo garantisce che la chiave API usata dal client Flutter (\texttt{anon\_key}) non abbia, di per sé, alcun permesso diretto di leggere o scrivere dati.

L'accesso ai dati è quindi interamente mediato da un'API sicura:
\begin{itemize}
    \item \textbf{Funzioni RPC \texttt{SECURITY DEFINER}}: per la maggior parte delle operazioni sul database.
    \item \textbf{Edge Functions}: per operazioni complesse e interazioni con servizi esterni.
\end{itemize}
Queste funzioni, agendo da ``gatekeeper'', contengono al loro interno tutta la logica di autorizzazione (es.\ ``l'utente è l'organizzatore di questo contest?''), verificando l'identità del chiamante tramite \texttt{auth.uid()} prima di eseguire qualsiasi operazione.

\subsection{L'Eccezione Necessaria: le Policy RLS per il Servizio Realtime}\label{subsec:l'eccezione-necessaria:-le-policy-rls-per-il-servizio-realtime}
Esiste un'eccezione mirata a questa strategia.
Il servizio \textbf{Supabase Realtime} basa il suo meccanismo di autorizzazione proprio sulle policy RLS della tabella.
Per permettere a un client di ``ascoltare'' le modifiche, è stato necessario definire policy di \texttt{SELECT} granulari, ma esclusivamente per le due tabelle esposte al servizio Realtime: \texttt{messages} e \texttt{voting\_sessions}.
Questo è stato un adattamento necessario per integrare una funzionalità chiave, mantenendo la massima sicurezza su tutte le altre operazioni e tabelle.

Il seguente listato mostra la policy implementata per la tabella \texttt{voting\_sessions}.

% Esempio di policy RLS per il servizio Realtime.
\begin{longlisting}
    \begin{minted}{postgresql}
    -- Abilita l'accesso in lettura alla tabella `voting_sessions`
    -- solo per gli utenti autorizzati a ricevere aggiornamenti in tempo reale.
    CREATE POLICY "Allow read access to organizers and involved jurors"
    ON public.voting_sessions
    FOR SELECT
    USING (
      -- Condizione 1: L'utente è l'organizzatore del contest a cui la sessione appartiene.
      (EXISTS (
        SELECT 1 FROM public.contests c
        WHERE c.id = voting_sessions.contest_id AND c.organizer_id = auth.uid()
      ))
      OR
      -- Condizione 2: L'utente è un giurato che partecipa a questa specifica sessione di voto.
      (EXISTS (
        SELECT 1 FROM public.voting_session_jurors vsj
        WHERE vsj.voting_session_id = voting_sessions.id AND vsj.juror_id = auth.uid()
      ))
    );
    \end{minted}
    \caption{Policy RLS di tipo SELECT per la tabella ``voting\_sessions'', necessaria per il servizio Realtime.}
    \label{lst:rls_realtime_policy}
\end{longlisting}

\subsection{Misure di Sicurezza Aggiuntive}\label{subsec:misure-di-sicurezza-aggiuntive}
Oltre al modello ibrido RLS/RPC, la sicurezza del sistema è rafforzata da altre misure:
\begin{itemize}
    \item \textbf{Protezione delle API Esterne}: le chiavi segrete per servizi come Resend e Google Places sono memorizzate in Supabase e accessibili solo dalle Edge Functions.
    \item \textbf{Protezione dei file}: nessun file nello Storage è pubblico.
    L'accesso avviene solo tramite \emph{signed URL} (URL firmati), ovvero link \textbf{a tempo limitato e con permessi specifici}, generati dinamicamente lato server per ogni singola richiesta.
\end{itemize}

\subsubsection{Gestione dei Giurati Semplici: il pattern dell'autenticazione anonima}
Una delle sfide più complesse è stata quella di permettere a un \emph{Simple Juror} di votare in modo sicuro senza obbligarlo a creare un account permanente.
Risolvere questo problema solo con un token temporaneo sarebbe stato insicuro, poiché non avrebbe permesso di tracciare in modo affidabile se l'utente avesse già votato.

Per risolvere questa sfida, è stato implementato un pattern basato sull'\textbf{autenticazione anonima} offerta da Supabase.
Questo approccio è un'applicazione diretta del Principio del Minimo Privilegio.
Il flusso è il seguente:

\begin{enumerate}
    \item \textbf{Creazione dell'account anonimo}: Nella pagina di accesso, l'utente seleziona l'opzione per partecipare come giurato semplice.
    L'applicazione gli chiede di inserire il proprio nome e cognome.
    Subito dopo, il client esegue una chiamata a Supabase per effettuare un \textbf{sign-in anonimo}.

    \item \textbf{Autenticazione e Sessione JWT}: Supabase Auth crea un account utente a tutti gli effetti, ma marcato come ``anonimo'', e restituisce un JWT valido.
    Da questo momento, anche l'utente anonimo ha un \texttt{auth.uid()} univoco e una sessione autenticata.
    L'\texttt{AuthBloc} nel client entra in uno stato ``autenticato ma anonimo''.

    \item \textbf{Accesso alla Home Page Limitata}: L'utente viene reindirizzato a una home page specifica per giurati semplici.
    L'\texttt{AuthGuard}, riconoscendo lo stato ``anonimo'', applica una logica di navigazione restrittiva: permette l'accesso solo a questa home page e alle pagine strettamente necessarie per la votazione, bloccando tutte le altre sezioni dell'app (es.\ impostazioni, gestione contest).

    \item \textbf{Accesso alla Sessione di Voto}: Dalla sua home page, il giurato semplice ha due opzioni per accedere alla sessione di voto: può \textbf{scannerizzare un QR code} fornito dall'organizzatore o \textbf{inserire manualmente il token} della giuria.

    \item \textbf{Validazione e Votazione}: Una volta inserito il token (tramite QR o manualmente), la RPC \texttt{juror\_access\_voting\_as\_simple\_juror} lo valida.
    Se il token è corretto, l'utente viene reindirizzato alla procedura di voto, dove può esprimere le sue preferenze.

    \item \textbf{Cancellazione dell'account al logout}: Al termine della procedura di voto, l'utente viene reindirizzato nuovamente alla sua home page limitata.
    Se a questo punto effettua il \emph{logout}, viene invocata la Edge Function \texttt{user-delete-account}.
    Poiché l'utente possiede un JWT valido, è autorizzato a chiamare questa funzione, che elimina in modo sicuro e permanente l'account anonimo appena utilizzato, garantendo che nessun dato temporaneo rimanga nel sistema.
\end{enumerate}

Questo pattern, oltre a garantire un'esperienza utente fluida e un alto livello di sicurezza, è stato progettato con un'importante considerazione sulla \textbf{scalabilità futura}.
La gestione del giurato semplice tramite un vero account, seppur anonimo, pone le basi per future evoluzioni, come l'implementazione di una funzionalità che permetta all'utente di ``reclamare'' il proprio account anonimo e \textbf{convertirlo in un account permanente}.

I vantaggi principali di questo approccio sono:
\begin{itemize}
    \item \textbf{Esperienza Utente Fluida}: L'utente non percepisce la creazione di un account, ma solo l'inserimento del proprio nome, seguito da un'interfaccia chiara per accedere alla votazione.
    \item \textbf{Sicurezza Robusta}: Ogni azione del giurato semplice è autenticata e legata a un \texttt{uid} univoco, prevenendo abusi.
    \item \textbf{Riutilizzo e Scalabilità dell'architettura}: Permette di riutilizzare la stessa logica di sicurezza (RLS, RPC, \texttt{AuthGuard}) sia per gli utenti normali che per quelli anonimi e apre la strada a future evoluzioni del sistema di account.
\end{itemize}

\section{Progettazione e Implementazione degli Strumenti di Amministrazione}
\label{sec:admin_tools}

Oltre all'applicazione destinata agli utenti finali, un sistema software robusto richiede strumenti adeguati per la sua gestione, moderazione e manutenzione.
In linea con l'obiettivo di concentrare le risorse di sviluppo sull'applicazione principale destinata agli utenti finali, per la gestione del sistema si è adottata una \textbf{strategia di sviluppo rapido}: invece di investire tempo nella creazione di un pannello di amministrazione da zero, si è scelto di utilizzare \textbf{Retool}, una piattaforma low-code specializzata nella creazione rapida di strumenti interni.

\subsection{Esigenza di un Pannello di Amministrazione}\label{subsec:esigenza-di-un-pannello-di-amministrazione}
Sebbene la dashboard di Supabase offra un accesso completo al database, essa è uno strumento pensato per gli sviluppatori, non per gli amministratori di sistema.
La necessità di un pannello di amministrazione dedicato nasce quindi da requisiti operativi critici che, pur essendo tecnicamente possibili, \textbf{non sono gestibili in modo semplice, sicuro e rapido} tramite l'interfaccia standard di Supabase:
\begin{itemize}
    \item \textbf{Gestione Utenti}: la capacità di visualizzare tutti gli utenti registrati, verificarne l'attività e, in caso di abuso, procedere con l'eliminazione dell'account.
    \item \textbf{Moderazione Contenuti}: la possibilità per un amministratore di ispezionare e, se necessario, rimuovere forzatamente un contest che viola i termini di servizio.
\end{itemize}

\subsection{Scelta di Retool per lo Sviluppo Rapido}\label{subsec:scelta-di-retool-per-lo-sviluppo-rapido}
La scelta di Retool è stata strategica e motivata da:
\begin{itemize}
    \item \textbf{Velocità di Sviluppo}: ha permesso di costruire un'interfaccia admin funzionale in poche ore, invece che in settimane, concentrando le risorse di sviluppo sull'applicazione principale.
    \item \textbf{Integrazione Nativa}: si integra nativamente con database PostgreSQL e API REST, permettendo di interfacciarsi facilmente con l'infrastruttura Supabase.
    \item \textbf{Sicurezza Configurabile}: offre meccanismi per gestire segreti (come le chiavi API) e configurare le chiamate in modo sicuro.
\end{itemize}

\subsection{Architettura di Comunicazione Sicura tra Retool e Supabase}\label{subsec:architettura-di-comunicazione-sicura-tra-retool-e-supabase}
La sicurezza è stata la priorità nella progettazione dell'integrazione.
L'accesso ai dati non è mai diretto, ma mediato da un'API sicura esposta dal backend Supabase, con due meccanismi distinti:

\subsubsection{Accesso tramite Utente Dedicato e Funzioni RPC}
Per le operazioni che richiedono solo l'accesso al database (es.\ visualizzare la lista di utenti), è stato creato un utente PostgreSQL dedicato, \texttt{retool}, con permessi minimi.
Le query da Retool invocano funzioni RPC con prefisso \texttt{admin\_*}, che al loro interno verificano che il chiamante (\texttt{session\_user}) sia esattamente \texttt{retool}, bloccando qualsiasi altro tentativo di accesso.

\begin{figure}[H]
    \centering
    \includegraphics[width=0.8\textwidth]{capitoli/figure/diagrams/retool-rpc-call}
    \caption{Flusso di una chiamata sicura da Retool a una Funzione RPC.}
    \label{fig:retool_rpc_call}
\end{figure}

% Inserire qui un estratto di codice PL/pgSQL che mostra la validazione del session_user in una funzione RPC amministrativa.
\begin{longlisting}
    \begin{minted}{postgresql}
    -- Estratto da una funzione come 'admin_get_all_accounts()'

    CREATE OR REPLACE FUNCTION admin_get_all_accounts()
    RETURNS SETOF profiles -- o un altro tipo di ritorno
    LANGUAGE plpgsql
    SECURITY DEFINER
    SET search_path = public
    AS $$
    BEGIN
      -- --- 1. CONTROLLO DI SICUREZZA ---
      -- Questa è la prima operazione eseguita dalla funzione.
      -- 'session_user' è una variabile speciale di PostgreSQL che contiene
      -- il nome dell'utente della connessione corrente (in questo caso, 'retool').

      IF session_user <> 'retool' THEN
        -- Se il chiamante non è l'utente 'retool' designato,
        -- la funzione solleva un'eccezione e interrompe immediatamente l'esecuzione.
        RAISE EXCEPTION 'Permission denied: This function can only be called by the admin user.';
      END IF;

      -- --- 2. ESECUZIONE DELLA LOGICA ---
      -- Se il controllo di sicurezza passa, la funzione procede
      -- con l'esecuzione della query per recuperare i dati.
      RETURN QUERY
        SELECT * FROM public.profiles;

    END;
    $$;
    \end{minted}
    \caption{Validazione del chiamante all'interno di una funzione RPC per l'amministrazione.}
    \label{lst:rpc_admin_validation}
\end{longlisting}

\subsubsection{Accesso tramite Edge Functions e API Key}
Per le operazioni che richiedono un'orchestrazione complessa o l'interazione con servizi esterni al database, come l'invio di email o la manipolazione di file nello Storage, sono state implementate delle \textbf{Edge Functions}.
Funzioni critiche come \texttt{admin-delete-account} e \texttt{admin-delete-contest} rientrano in questa categoria, poiché devono non solo modificare i record del database, ma anche eliminare tutti i file associati (immagini, file) dallo Storage per garantire una pulizia completa.

Queste funzioni sono protette da un meccanismo di autenticazione server-to-server per garantire che possano essere invocate solo da un client autorizzato come Retool:
\begin{enumerate}
    \item Retool, nell'eseguire l'azione, invia una richiesta HTTP \texttt{POST} alla Edge Function, includendo un header personalizzato: \texttt{X-API-Key}.
    \item Il valore di questa chiave è un segreto memorizzato sia in Retool che in Supabase.
    \item La Edge Function, come primo passo, confronta la chiave ricevuta con quella memorizzata nei suoi segreti.
    Se non corrispondono, la richiesta viene immediatamente respinta con un errore 401 (Unauthorized).
\end{enumerate}
Questo approccio garantisce che solo un'applicazione autorizzata possa invocare queste funzioni, proteggendo il sistema da accessi impropri.

\begin{figure}[H]
    \centering
    \includegraphics[width=0.8\textwidth]{capitoli/figure/diagrams/retool-edge-call}
    \caption{Flusso di una chiamata sicura da Retool a una Edge Function, protetta da una chiave API.}
    \label{fig:retool_edge_call}
\end{figure}

% Inserire qui un estratto di codice TypeScript che mostra la validazione della X-API-Key in una Edge Function amministrativa.
\begin{longlisting}
    \begin{minted}{typescript}
    // Estratto da una Edge Function come 'admin-delete-account/index.ts'

	// ... (all'interno della funzione Deno.serve)
	try {
		// --- 1. API KEY AUTHORIZATION ---
		// Questo pattern è usato per tutte le chiamate server-to-server da Retool.

		// Legge la chiave API inviata dal client (Retool) nell'header.
		const apiKey = req.headers.get('X-API-Key');

		// Recupera la chiave segreta memorizzata in modo sicuro nel Vault di Supabase.
		const secretKey = Deno.env.get('RETOOL_API_KEY');

		// Primo controllo: verifica che il segreto sia stato configurato.
		if (!secretKey) {
			throw new Error("RETOOL_API_KEY is not set in Supabase secrets.");
		}

		// Secondo controllo: confronta la chiave ricevuta con il segreto.
		// Se non corrispondono, la richiesta viene immediatamente respinta.
		if (apiKey !== secretKey) {
			return new Response(
			JSON.stringify({ error: "Unauthorized. Invalid API Key." }),
			{
				headers: { ...corsHeaders, "Content-Type": "application/json" },
				status: 401, // 401 Unauthorized
			}
			);
		}

		// --- 2. PROSECUZIONE DELLA LOGICA ---
		// Se la validazione ha successo, la funzione procede con
		// l'esecuzione delle operazioni amministrative.

		// ... (logica per eliminare l'utente, i file, ecc.)

	} catch (error) {
		// ... (gestione degli errori)
	}
    \end{minted}
    \caption{Validazione della chiave API all'interno di una Edge Function amministrativa.}
    \label{lst:api_key_validation}
\end{longlisting}


\subsection{Funzionalità Implementate nel Pannello}\label{subsec:funzionalita-implementate-nel-pannello}
La dashboard sviluppata in Retool fornisce all'amministratore le seguenti funzionalità chiave:
\begin{itemize}
    \item \textbf{Una tabella} per la visualizzazione e la ricerca di tutti gli utenti registrati.
    \item \textbf{Una vista di dettaglio} per un singolo utente, che mostra i contest a cui è associato e i suoi ruoli.
    \item \textbf{Un pulsante ``Elimina Utente''}, che invoca in modo sicuro la Edge Function \texttt{admin-delete-account}.
    \item \textbf{Una tabella di tutti i contest}, con un pulsante ``Elimina Contest'' che invoca la Edge Function \texttt{admin-delete-contest}.
\end{itemize}